\documentclass[stu,12pt,floatsintext]{apa7}

\usepackage[utf8]{inputenc}
\usepackage[T1]{fontenc}
\usepackage[spanish,es-noquoting]{babel}
\usepackage{amsmath,amssymb,amsfonts}
\usepackage{graphicx}
\usepackage{booktabs}
\usepackage{csquotes}
\usepackage{url}

\usepackage[style=apa,backend=biber]{biblatex}
\DeclareLanguageMapping{spanish}{spanish-apa}
\addbibresource{referencias.bib}

\graphicspath{{figuras/}}

% Las figuras las genera src/experiments/ablation.py desde los CSV. Mientras
% no existan, el documento debe seguir compilando: en su lugar se dibuja un
% marcador con el nombre del archivo que falta.
\newcommand{\figurascoal}[2]{%
  \IfFileExists{figuras/#1}%
    {\includegraphics[width=#2]{#1}}%
    {\fbox{\parbox[c][3cm][c]{#2}{\centering\ttfamily
       figura pendiente\\\detokenize{#1}}}}%
}

% --- Metadatos de portada (clase apa7, opcion stu) ---
\title{Clasificacion de objetos mediante algebra lineal:
       estudio comparativo de tres representaciones de caracteristicas
       bajo un clasificador de minimos cuadrados}
\shorttitle{Clasificacion de objetos mediante algebra lineal}
\author{Jhonnatan}
\affiliation{Facultad de Ingenieria, Universidad Cooperativa de Colombia}
\course{Algebra Lineal}
% TODO: reemplazar por el nombre completo del docente.
\professor{Nombre del docente}
% TODO: reemplazar por la fecha de entrega.
\duedate{Fecha de entrega}

\begin{document}

\maketitle                              % portada APA 7 (estudiante)
\newpage
\thispagestyle{empty}
\begin{center}

{\large\textbf{Clasificacion de objetos mediante algebra lineal:
estudio comparativo de tres representaciones de caracteristicas
bajo un clasificador de minimos cuadrados}}

\vspace{2cm}

Jhonnatan

\vspace{2cm}

Trabajo presentado como requisito para la asignatura de Algebra Lineal

\vspace{2cm}

Docente\\
Nombre del docente          % TODO: nombre completo del docente

\vspace{2cm}

Universidad Cooperativa de Colombia\\
Facultad de Ingenieria\\
Ingenieria de Sistemas\\
Villavicencio, Meta\\
2026                        % TODO: verificar el ano de entrega

\end{center}
\newpage


\begin{abstract}
% TODO: reemplazar este parrafo por el resumen real, de 150 a 250 palabras en
% espanol. Debe cubrir: problema, metodo, resultado principal cuantificado y
% conclusion. Los valores numericos salen de docs/tablas/resultados_via.tex,
% que a su vez se genera desde results/ablacion_via.csv: no escribir cifras a
% mano. El entorno abstract no puede quedar vacio o la clase apa7 falla.
Resumen pendiente de redaccion.
\end{abstract}

\keywords{algebra lineal, clasificacion de objetos,
analisis de componentes principales, minimos cuadrados,
descomposicion en valores singulares, vision por computador}

\newpage
\section*{Abstract}
% TODO: version en ingles del resumen.

\noindent\textit{Keywords:} linear algebra, object classification,
principal component analysis, least squares, singular value
decomposition, computer vision.
\newpage

\section{Introduccion}

% TODO: 3 a 5 parrafos. Guion sugerido:
%   1. Contexto: clasificacion automatica de objetos en linea de produccion.
%   2. Problema: cuanto del desempeno viene del clasificador y cuanto de la
%      representacion de entrada.
%   3. Antecedentes: eigenfaces \parencite{turk1991}, minimos cuadrados como
%      clasificador \parencite{boyd2018}, datasets \parencite{nene1996,muresan2018}.
%   4. Aporte: mismo clasificador lineal sobre tres representaciones, con el
%      motor implementado en NumPy sin bibliotecas de aprendizaje automatico.
%   5. Organizacion del documento.

% Ejemplo de cita para conservar el estilo APA autor-fecha:
% El metodo de eigen-objetos se deriva del trabajo de \textcite{turk1991}.

\section{Objetivos}

\subsection{Objetivo general}

Desarrollar un sistema de clasificacion de objetos por tamano, color, forma e
identidad, implementado mediante operaciones de algebra lineal, que permita
evaluar comparativamente el aporte de tres representaciones de caracteristicas
distintas bajo un mismo clasificador de minimos cuadrados regularizado.

\subsection{Objetivos especificos}

\begin{enumerate}
\item Implementar tres extractores de caracteristicas ---descriptores
geometricos y cromaticos, proyeccion sobre eigen-objetos mediante
descomposicion en valores singulares, y representaciones de una red
convolucional preentrenada--- empleando unicamente operaciones matriciales
programadas sobre NumPy.

\item Construir un clasificador multiclase de minimos cuadrados con
regularizacion de Tikhonov, verificando su equivalencia numerica con una
implementacion de referencia y caracterizando el efecto del parametro de
regularizacion sobre el numero de condicion de la matriz normal.

\item Evaluar el desempeno comparativo de las tres representaciones mediante
un estudio de ablacion con validacion cruzada, cuantificando la contribucion
de la representacion frente a la del clasificador en la exactitud final del
sistema.
\end{enumerate}

\section{Marco teorico}

% TODO: desarrollar cada subseccion. Las formulas ya estan planteadas y
% coinciden con las docstrings del codigo fuente; el texto debe explicarlas.

\subsection{La imagen como objeto algebraico}

Una imagen digital en color es un tensor $I \in \mathbb{R}^{H \times W \times 3}$.
% TODO: desarrollar.

\subsection{Transformaciones lineales de color}

La conversion de RGB a YCbCr, en la recomendacion ITU-R BT.601, es una
transformacion afin cuya parte lineal es una matriz de $3 \times 3$:

\begin{equation}
\begin{bmatrix} Y \\ C_b \\ C_r \end{bmatrix}
=
\begin{bmatrix}
 0.2990 &  0.5870 &  0.1140 \\
-0.1687 & -0.3313 &  0.5000 \\
 0.5000 & -0.4187 & -0.0813
\end{bmatrix}
\begin{bmatrix} R \\ G \\ B \end{bmatrix}
+
\begin{bmatrix} 0 \\ 128 \\ 128 \end{bmatrix}.
\end{equation}

La conversion de RGB a HSV, en cambio, \emph{no} es lineal: el tono y la
saturacion se definen a partir de $\max(R,G,B)$, $\min(R,G,B)$ y cocientes
entre esas cantidades, operaciones que no admiten representacion matricial.
Esa es la razon por la que la segmentacion de este trabajo opera en YCbCr.
% TODO: ampliar.

\subsection{Momentos y matriz de inercia}

Con los momentos centrales de segundo orden de la mascara binaria se construye

\begin{equation}
M = \frac{1}{\mu_{00}}
\begin{bmatrix} \mu_{20} & \mu_{11} \\ \mu_{11} & \mu_{02} \end{bmatrix},
\end{equation}

matriz simetrica y semidefinida positiva cuyos autovectores son las direcciones
principales del objeto y cuyos autovalores miden la dispersion a lo largo de
esas direcciones.
% TODO: desarrollar.

\subsection{Descomposicion en valores singulares y eigen-objetos}

Sea $A \in \mathbb{R}^{m \times d}$ la matriz de imagenes centradas. Su
descomposicion en valores singulares $A = U \Sigma V^{T}$ entrega en las
columnas de $V$ los autovectores de la matriz de covarianza
$C = A^{T}A/(m-1)$. Cuando $m \ll d$ se aplica la identidad de
\textcite{turk1991}: si $A A^{T} u = \mu\, u$, entonces

\begin{equation}
A^{T}A \,(A^{T}u) = \mu \,(A^{T}u),
\end{equation}

de modo que $v = A^{T}u / \lVert A^{T}u \rVert$ es autovector de la matriz
grande sin necesidad de construirla.
% TODO: desarrollar.

\subsection{Minimos cuadrados regularizados}

El clasificador resuelve

\begin{equation}
\min_{W} \; \lVert XW - Y \rVert_{F}^{2} + \lambda \lVert W \rVert_{F}^{2},
\qquad
W = (X^{T}X + \lambda I)^{-1} X^{T} Y ,
\end{equation}

que con $\lambda = 0$ coincide con $W = X^{+}Y$, donde $X^{+}$ es la
pseudoinversa de Moore-Penrose \parencite{boyd2018}.
% TODO: desarrollar, incluir la interpretacion geometrica de la proyeccion
% sobre el espacio columna de X y la frontera de decision como hiperplano.

\subsection{Numero de condicion y regularizacion}

\begin{equation}
\operatorname{cond}_2(X^{T}X + \lambda I)
= \frac{\sigma_{\max} + \lambda}{\sigma_{\min} + \lambda}
\end{equation}

decrece monotonamente con $\lambda$: la regularizacion desplaza todo el
espectro y estabiliza la solucion a costa de sesgo.
% TODO: desarrollar.

\subsection{Mapeo de conceptos}

\begin{table}[htbp]
\caption{Conceptos de algebra lineal y su punto de aplicacion en el sistema}
\label{tab:mapeo-conceptos}
\begin{tabular}{ll}
\toprule
Concepto & Donde se aplica \\
\midrule
Matrices y tensores & Imagen como $I \in \mathbb{R}^{H \times W \times 3}$ \\
Producto matriz-vector & Conversion RGB a YCbCr y RGB a gris \\
Autovalores y autovectores & Matriz de inercia $2\times2$, ejes principales \\
Matriz de covarianza & Estandarizacion y analisis de componentes principales \\
Descomposicion en valores singulares & Eigen-objetos, reduccion de 4096 a $k$ \\
Relacion entre $A^{T}A$ y $AA^{T}$ & Truco de Turk-Pentland para $m \ll n$ \\
Espacios y subespacios vectoriales & Espacio de caracteristicas, subespacio propio \\
Proyeccion ortogonal & Representacion en el subespacio de eigen-objetos \\
Minimos cuadrados & Entrenamiento del clasificador \\
Pseudoinversa de Moore-Penrose & Solucion cerrada $W = X^{+}Y$ \\
Rango e independencia lineal & Deteccion de caracteristicas redundantes \\
Numero de condicion & Diagnostico de $X^{T}X$ mal condicionada \\
Regularizacion de Tikhonov & Estabilizacion mediante $+\lambda I$ \\
Norma euclidiana & Baseline k-NN, distancia en el espacio propio \\
Hiperplanos separadores & Frontera de decision del clasificador lineal \\
\bottomrule
\end{tabular}
\note{Elaboracion propia.}
\end{table}


\section{Metodologia}

% TODO: redactar en pasado y en tercera persona, siguiendo el orden del
% pipeline implementado.

\subsection{Datos}

% TODO: COIL-100 \parencite{nene1996} y Fruits-360 \parencite{muresan2018}.
% Justificar el subconjunto de diez clases con solapamiento deliberado y el
% balanceo a trescientas imagenes por clase.

\subsection{Segmentacion y normalizacion}

% TODO: umbral sobre la distancia al color de fondo en YCbCr, operaciones
% morfologicas, contorno de area maxima, recorte cuadrado con fondo neutro.

\subsection{Representaciones evaluadas}

% TODO: via A (dieciseis descriptores geometricos y cromaticos), via B
% (eigen-objetos sobre imagenes de 64x64), via C (embeddings de
% MobileNetV3-Small congelado \parencite{howard2019} reducidos por PCA propia).

\subsection{Clasificador}

% TODO: minimos cuadrados multiclase con regularizacion de Tikhonov;
% estandarizacion ajustada unicamente con el conjunto de entrenamiento.

\subsection{Protocolo experimental}

% TODO: particion estratificada con semilla fija, validacion cruzada de cinco
% pliegues, tres barridos de ablacion y baselines de comparacion.

\subsection{Implementacion}

% TODO: NumPy como unico motor; sklearn empleado solo como referencia de
% verificacion numerica; torch empleado solo como extractor congelado.

\section{Resultados}

% Las cifras de este capitulo se inyectan desde los CSV del estudio de
% ablacion, mediante las tablas generadas en docs/tablas/. No escribir valores
% numericos a mano: regenerar con
%     python -m src.experiments.ablation
% y volver a compilar.

\subsection{Comparacion entre representaciones}

\begin{table}[htbp]
\caption{Exactitud por representacion bajo validacion cruzada}
\label{tab:resultados-via}
\begin{tabular}{llrrr}
\toprule
Via & Representacion & $M$ exactitud & $DE$ & $M$ macro-F1 \\
\midrule
A & Fisicas & 0.841 & 0.015 & 0.841 \\
B & Eigen & 0.900 & 0.011 & 0.899 \\
\bottomrule
\end{tabular}
\note{Validacion cruzada estratificada de 5 pliegues sobre el mismo
clasificador de minimos cuadrados con $\lambda = 0.01$ y $k = 50$.
Elaboracion propia.}
\end{table}


% TODO: describir la tabla en prosa.

\begin{figure}[htbp]
\caption{Exactitud por representacion bajo validacion cruzada}
\figurascoal{ablacion_via.png}{0.6\textwidth}
\note{Media y desviacion estandar sobre cinco pliegues estratificados.
Elaboracion propia.}
\end{figure}

\subsection{Efecto del numero de componentes}

\begin{figure}[htbp]
\caption{Exactitud y varianza explicada en funcion de $k$}
\figurascoal{ablacion_k.png}{\textwidth}
\note{Vias B y C. La linea discontinua marca el noventa por ciento de
varianza explicada acumulada. Elaboracion propia.}
\end{figure}

% TODO: comentar el punto de saturacion.

\subsection{Efecto de la regularizacion}

\begin{figure}[htbp]
\caption{Exactitud y numero de condicion en funcion de $\lambda$}
\figurascoal{ablacion_lambda.png}{\textwidth}
\note{El panel derecho muestra $\operatorname{cond}_2(X^{T}X+\lambda I)$ en
escala logaritmica. Elaboracion propia.}
\end{figure}

% TODO: comentar el optimo y la caida del numero de condicion.

\subsection{Comparacion contra los baselines}

\begin{table}[htbp]
\caption{Comparacion contra los baselines}
\label{tab:baselines}
\begin{tabular}{llr}
\toprule
Baseline & Via & Exactitud \\
\midrule
aleatorio & - & 0.100 \\
knn\_k5 & A & 0.927 \\
minimos\_cuadrados & A & 0.821 \\
knn\_k5 & B & 0.976 \\
minimos\_cuadrados & B & 0.872 \\
\bottomrule
\end{tabular}
\note{Todos los valores sobre la misma particion de prueba.
Elaboracion propia.}
\end{table}


% TODO: comentar la brecha entre la via C y la cabeza original de MobileNetV3.

\subsection{Analisis cualitativo}

\begin{figure}[htbp]
\caption{Primeros dieciseis eigen-objetos}
\figurascoal{eigenobjetos.png}{0.6\textwidth}
\note{Cada panel indica la fraccion de varianza explicada. Elaboracion propia.}
\end{figure}

\begin{figure}[htbp]
\caption{Proyeccion sobre las dos primeras componentes principales}
\figurascoal{scatter_pca.png}{0.6\textwidth}
\note{Color por clase. Elaboracion propia.}
\end{figure}

\begin{figure}[htbp]
\caption{Matriz de confusion de la via A}
\figurascoal{confusion_A.png}{0.55\textwidth}
\note{Filas normalizadas por clase real. Elaboracion propia.}
\end{figure}

\begin{figure}[htbp]
\caption{Matriz de confusion de la via B}
\figurascoal{confusion_B.png}{0.55\textwidth}
\note{Filas normalizadas por clase real. Elaboracion propia.}
\end{figure}

\begin{figure}[htbp]
\caption{Matriz de confusion de la via C}
\figurascoal{confusion_C.png}{0.55\textwidth}
\note{Filas normalizadas por clase real. Elaboracion propia.}
\end{figure}

\section{Discusion}

% TODO: responder la pregunta de investigacion con los numeros obtenidos.
% Puntos que el trabajo debe discutir:
%   - Cuanto del desempeno depende de la representacion y cuanto del
%     clasificador, comparando las tres vias bajo el mismo clasificador.
%   - Por que el clasificador lineal alcanza, o no, el techo del extractor
%     congelado frente a la cabeza original de la red.
%   - La discrepancia entre vias en la demostracion en vivo como evidencia de
%     desplazamiento de dominio, no como defecto de implementacion: los
%     datasets tienen fondo controlado y la camara no.
%   - Limitaciones: fondo uniforme, un solo objeto por escena, ausencia de la
%     medicion de masa.

\section{Conclusiones}

% TODO: una conclusion por objetivo especifico, en el mismo orden en que
% fueron planteados, cada una respaldada por un resultado cuantificado del
% capitulo anterior.

\section{Trabajo futuro}

% TODO: ampliar cada linea en prosa.

\begin{itemize}
\item Incorporacion del peso real mediante celda de carga HX711 sobre ESP32.
La arquitectura ya contempla esa componente del vector de caracteristicas,
desactivada mientras no existan mediciones de masa.
\item Ajuste fino del extractor sobre imagenes del dominio propio, para cerrar
la brecha observada entre dataset y camara.
\item Banda transportadora fisica, con disparo por fotocelda en lugar del
disparador de estabilidad por software.
\item Comparacion contra una maquina de vectores de soporte con nucleo, para
medir cuanto aporta la no linealidad situada en el clasificador y no en la
representacion.
\end{itemize}


\printbibliography[title={Referencias}]

\end{document}
